% =====================================================================
% paper/sections/conclusion.tex
% §8 Conclusion（rev.1, 2026-05-01）
%
% 写作纪律（outline §10）：
%   - 4 句话收尾，几乎逐字拷自 outline §0.3 narrative spine
%   - 不报新数字、不引新文献、不开新论点
%   - 总长 ≤ 1/4 column
%
% 锚点：sec:conclusion（被 intro.tex 1 处 forward-ref）
% =====================================================================

\section{Conclusion}\label{sec:conclusion}

我们提出 ReBRAC-Q —— 一个 Q-normalized dual-penalty TD3+BC 变体 —— 用于水下 AUV wake field 导航这一 sim2real 任务。在 \texttt{auv\_nav} 环境（REMUS-100 6-DOF + Kármán 尾迹场 + s0 单点 DVL）下，ReBRAC-Q 的 deployable-only 协议性能与 privileged-critic 协议在统计意义下不可区分（\texttt{worldcomp-1000} 上 Welch's $p = 0.92$），且在 \texttt{crosscomp} 上相对 vanilla TD3+BC 提升 $+23.0$/$+32.2$pp。机制上，actor-side BC penalty 主导 mean uplift，critic-side BC penalty 不抬 mean 但通过 target-Q bootstrap 在两个 $\hat{Q}$ 符号相反的 dataset 上同时提供\emph{方向一致}的 OOD 抑制（dataset-invariant 稳定性），并对 outlier seed 提供独立救援；critic LayerNorm 与 dual penalty 是两个相互正交、抑制不同 failure mode 的稳定器。这一结果意味着，在 deployable sensor 严格受限的水下 AUV 部署设置下，离线 RL 不依赖任何 simulator-only privileged 通道也能把策略性能逼近 privileged-critic 上界 —— 在我们的实验 scope（sub-critical regime + efficiency-v2 reward；\S\ref{sec:limitations}）下，部署侧的策略训练\emph{不需要}知道壳积分等效水流；该结论的有效边界与 critical-regime 上的反例见 \S\ref{sec:limitations} L12 与 \S\ref{subsec:discussion_regime_boundary}。
